% ===========================
% Pustaka dan Format Umum
% ===========================

\usepackage{ragged2e}   % buat kontrol perataan teks
\usepackage{microtype}
\usepackage[a4paper, margin=1in]{geometry}  % Margin halaman A4
\usepackage{graphicx}                       % Menyisipkan gambar
\usepackage{tocbibind}                      % Tambah daftar pustaka ke TOC

% ===========================
% Font
% ===========================

\usepackage{fontspec}

\setmainfont[
  Path = conf/,
  Extension = .ttf,
  UprightFont = {times},
  BoldFont = {bold}
]{TimesNewRoman}

% ===========================
% Package Tambahan
% ===========================

\usepackage{float}
\usepackage{titlesec}
\usepackage{array}
\usepackage{parskip}
\usepackage{indentfirst}
\usepackage{setspace}
\usepackage{xcolor}

% ===========================
% Hyperref
% ===========================

\usepackage[
    colorlinks=true,
    linkcolor=blue,
    urlcolor=blue,
    citecolor=blue
]{hyperref}

% ===========================
% Bibliography
% ===========================

\usepackage[
    style=numeric,
    sorting=none,
    backend=biber
]{biblatex}

\addbibresource{Referensi.bib}

% ===========================
% Konfigurasi Tambahan
% ===========================

\AtBeginDocument{\justifying}

\setstretch{1.5}
\setlength{\parindent}{2em}
\setlength{\parskip}{0pt}

\hypersetup{hidelinks}

% ===========================
% Format Judul Bab/Subbab
% ===========================

\titleformat{\section}
{\centering\Large\bfseries}
{}{0em}{}

\titleformat{\subsection}
{\large\bfseries}
{}{0em}{}